%% ----------------------------------------------------------------
%% Introduction.tex
%% ---------------------------------------------------------------- 
\chapter{Introduction} \label{Chapter:Introduction}

% The recent development of low altitude hypersonic vehicle technology for both civilian and military applications has prompted a rejuvenated scientific interest in the field. The hypersonic environment is conventionally considered as flows above Mach 5, wherein flight vehicles encounter strong shock-wave compressions which cause an abrupt increase in the gas flow temperature. This gives rise to a \textit{`high-enthalpy'} environment; where the molecular dissociation in the atmosphere causes the flow to be in chemical and vibrational non-equilibrium. The intricate nature of high-enthalpy flows induce complex interactions between the flight vehicle's surface and the flow around, where the dissipation of large amount of the fluid kinetic energy into thermal energy within the viscous boundary layer. This consequently causes an acute rise in the peak temperature at small wall normal distances causing high temperature gradients in the boundary layer, generating higher values of surface heat flux. 

The recent development of low-altitude hypersonic vehicle technology for both civilian and military applications has prompted renewed scientific interest in the field. The hypersonic regime is conventionally defined as flows at speeds above Mach 5, where flight vehicles encounter strong shock-wave compressions that cause an abrupt increase in gas temperature. This gives rise to a \textit{high-enthalpy} environment, in which molecular dissociation in the atmosphere leads to chemical and vibrational nonequilibrium \citep{Candler2019}. The intricate nature of high-enthalpy flows induces complex interactions between the vehicle surface and the surrounding flow, characterised by the dissipation of a large amount of fluid kinetic energy into thermal energy within the viscous boundary layer. This, in turn, causes an acute rise in peak temperature over small wall-normal distances, leading to steep temperature gradients and elevated surface heat fluxes. 

From an engineering standpoint, the efficient design of such vehicles requires accurate prediction of the aerothermodynamic loads acting on the surface — quantities such as heat flux, pressure, and shear stress, from which the aerodynamic forces and moments must be estimated. These factors govern not only the vehicle’s aerodynamic performance but also dictate the selection and sizing of the \gls{TPS}, which must protect the structure from extreme thermal loads and steep temperature gradients. \fref{figure:trajectory} illustrates a comparison of the trajectories of various existing and experimental vehicles, including a representative \gls{HGV} \citep{Rizvi2015} and the Shuttle Orbiter. For each trajectory, the heat rate limit is determined by the durability and capability of the chosen \gls{TPS}. Throughout its journey, the \gls{HGV} experiences some Nitrogen dissociation, significant Oxygen dissociation, and pronounced vibrational excitation. These thermochemical processes significantly alter the physics of the boundary layer, affecting not only the temperature and velocity profiles but also the transport properties such as viscosity, thermal conductivity, and species diffusion rates. As a result, conventional modelling approaches based on perfect gas assumptions are inadequate, and a detailed treatment of the thermochemical state becomes necessary to accurately predict the aerothermodynamic loads on the vehicle surface.

% \begin{figure}[h!]
%     \centering
%     % \includegraphics[width=13cm]{resources/images/report2022tempranges.pdf}
%     \includegraphics[scale=0.34]{resources/figures/introduction/schematictrajectory.svg}
%     \caption{Schematic of a small-sized waverider going at hypersonic speed in the Earth’s atmosphere experiencing real gas phenomena.}
%     \label{figure:trajectory}
% \end{figure}


\begin{figure}[hbt!]
    \centering
    \fontsize{8}{9}\selectfont\includesvg[width=0.96\textwidth]{resources/figures/introduction/schematictrajectory.svg}
    \caption{Schematic of a small-sized waverider going at hypersonic speed in the Earth’s atmosphere experiencing real gas phenomena adapted from \citep{Anderson2019} with the HGV glide path imposed from \citep{Rizvi2015}.}
    \label{figure:trajectory}
\end{figure}


The flow conditions downstream of the leading-edge stagnation region pose significant challenges for the design of the \gls{TPS} and the vehicle’s structural integrity. At the leading edge, a strong shock compresses and heats the air, triggering nonequilibrium phenomena immediately behind the shock. Translational modes equilibrate rapidly with the local state, but vibrational and chemical relaxation occur over much longer timescales. As the flow moves along the surface, these relaxation processes remain incomplete, resulting in a non-uniform thermochemical state throughout the glide path.

At low altitudes and high speeds, the elevated atmospheric density results in high Reynolds numbers, making the flow highly susceptible to turbulence. Even small perturbations, such as \gls{TPS} ablation, receptivity in the stagnation region, or freestream disturbances, can trigger laminar–turbulent transition, substantially increasing the thermal loads on the vehicle’s surface. Predicting transition onset is inherently challenging and becomes even more complex under these conditions, where high stagnation enthalpies intensify real gas effects, including vibrational excitation and chemical nonequilibrium. These additional modes of energy exchange alter the boundary-layer stability characteristics, adding further uncertainty. Since turbulence can dramatically amplify aerothermodynamic loads and increase the risk of structural failure, a detailed understanding of the transition pathways and the mechanisms driving disturbance amplification is essential for robust thermal protection system design and overall vehicle survivability.


Capturing the correct physics under such extreme conditions presents significant challenges, both experimentally and computationally~\citep{DiRenzo2021}. Only a limited number of experimental facilities can replicate the high-energy environments needed for accurate representation, and this is further hindered by the substantial power demands of the test rigs and the requirement for quiet flow conditions in transition studies. Consequently, most research has focused primarily on the early stages of the transition process, leaving a comprehensive framework for understanding the full transition mechanism including the pathways through which instability modes breakdown to turbulence, study of which is still lacking in hypersonic boundary layers. A clear understanding of boundary layer behaviour at high Mach numbers and low altitudes, where elevated temperatures and turbulent regimes may develop, is crucial for ensuring safe, sustainable, and efficient vehicle design. These aspects remain largely unexplored, both individually and in combination, highlighting the need for comprehensive studies that address their coupled effects.

% The present work addresses this gap by using \gls{DNS}, in combination with \gls{LST}, to better identify the role of real gas effects, vibrational relaxation, and nonequilibrium phenomena, and to clarify the mechanisms and pathways by which disturbances amplify and breakdown to turbulence.


% The flow conditions downstream of the leading-edge stagnation region present significant challenges to the design of the \gls{TPS} and, more broadly, to the vehicle's structural integrity. At the leading edge, the vehicle encounters a strong shock that compresses and heats the surrounding air, producing a sudden rise in thermodynamic properties and triggering a cascade of nonequilibrium phenomena. Immediately behind the shock, the gas is driven far from equilibrium. While translational modes equilibrate rapidly and respond almost instantaneously to the local thermodynamic state, vibrational and chemical relaxation occur over much longer timescales. As the flow continues along the vehicle surface, these relaxation processes remain incomplete, resulting in a highly nonuniform thermochemical state during the glide path.

% Furthermore, at low altitudes and high speeds, the elevated atmospheric density leads to high Reynolds numbers, making the flow susceptible to turbulence. Under these conditions, small perturbations such as ablation or erosion of the \gls{TPS} material, stagnation region receptivity, or freestream disturbances can trigger the laminar-turbulent transition, substantially increasing thermal loads on the vehicle. Transitional flows in hypersonic regimes have been studied extensively since the mid-20th century, particularly in the context of re-entry vehicles, with continuing research aimed at understanding transition mechanisms and sensitivities. However, predicting transition onset at low altitudes, where freestream stagnation enthalpies are high, remains challenging and un explored. Since turbulence can significantly amplify aerothermodynamic loads and pose a risk of structural failure, a detailed understanding of transition processes is essential for robust and reliable vehicle design.



% To aid understanding of these phenomena, simplified geometries such as flat plates and wedges are commonly employed to study the fundamental behaviour of boundary layers and shock–boundary-layer interactions under high-enthalpy conditions.


% Amongst the critical considerations for vehicle design are the general shape of the vehicle, the shape of the nose, and the leading edges of the wings and other aerodynamic surfaces. Leading edges are of particular importance as they encounter the stagnation region, caused by the abrupt deceleration of the flow impinging on the body, leading to maximum temperature and pressure zones near the moving body. Generally, to reduce the overall drag on a vehicle in supersonic flows, a sharp, slender nose geometry generating a weak shock is preferred. However, in the hypersonic regime, a sharp geometry leads to an attached shock, which significantly increases the tip stagnation temperature — potentially beyond the melting point of the leading-edge material. To distribute this excessive thermal load, blunt nose shapes producing strong shock waves were historically adopted in high-altitude hypersonic vehicles. Although new advancements in material science have enabled the use of sharp noses and leading edges, infinite sharpness on an aircraft is practically unachievable due to manufacturing limitations, resulting in most structures possessing a certain degree of bluntness. Therefore, accurately understanding the flow behaviour around a blunt leading edge is of critical importance.


% \begin{figure}
%     \centering
%     \fontsize{5}{6}\selectfont\includesvg[width=0.71\textwidth]{resources/figures/introduction/schematictrajectory.svg}
%     \captionof{figure}{Schematic of a small-sized waverider going at hypersonic speed in the Earth's atmosphere.}
%     \label{fig:schematicflatplate}
% \end{figure}


% \newpage

\section{High enthalpy flows}
Vehicles operating in hypersonic flows encounter intense shock-wave compression, forming a high-enthalpy gas environment characterised by complex physicochemical processes. Under these extreme conditions, the coupling between macroscopic flow properties and the microscopic internal energy modes of the gas becomes significantly more pronounced, necessitating more advanced modelling than that used for classical, low-enthalpy flows \citep{Candler2019}. These internal energy modes — including translation, rotation, vibration, dissociation and exchange, electronic excitation, radiation, and ionisation — absorb, store, and release energy at discrete, quantised levels \cite[ch.~11]{Anderson2019}. At sufficiently high temperatures, both vibrational and chemical modes may lag behind the translational–rotational modes, and may also exhibit different relaxation times relative to one another, leading to a thermochemical nonequilibrium state. This delayed relaxation toward equilibrium introduces modelling challenges due to the intricate coupling between thermal and chemical processes.

% In hypersonic flows, these internal modes are predominantly endothermic, acting as heat sinks that absorb energy—contrary to exothermic flows such as combustion, where energy is released. The accessibility of each mode depends on the enthalpy associated with the vehicle's flight conditions \citep{Gupta1990}. Furthermore, each mode requires a threshold energy before it becomes active, a process often described as "mode activation". 

% Understanding the modelling implications first requires awareness of the physical processes that occur within gas species as temperature increases. In conventional flows (i.e., those operating below approximately 800 K), vibrational modes are largely inactive, and translational and rotational modes dominate, determining the total internal energy. Under these conditions, the gas can be treated as calorically perfect, with constant specific heats and a constant ratio of specific heats, $\gamma$, since $c_p$ and $c_v$ remain independent of temperature. Above 800 K, vibrational modes begin to activate, significantly altering the thermodynamic properties of the gas. In this regime, the gas behaves as thermally perfect—where $c_p$ and $c_v$ become functions of temperature—and calorically imperfect due to the variable specific heat ratio. With further temperature increases, chemical reactions are triggered, leading to the dissociation of atmospheric species and increased compositional complexity \citep{Fletcher,Scoggins2020}. This progression of vibrational and chemical excitation in Earth's atmosphere is illustrated in \fref{figure:tempranges}.

In hypersonic flows, internal energy modes are predominantly endothermic, acting as heat sinks that absorb energy unlike exothermic flows such as combustion, where energy is released. The accessibility of each mode depends on the enthalpy associated with the vehicle's flight conditions~\citep{Gupta1990}, with each mode requiring a threshold energy before becoming active. As temperature increases, the internal behaviour of gas species changes significantly; a detailed depiction of this behaviour in the Earth's atmosphere is illustrated in \fref{figure:tempranges}. In conventional flows (i.e., below approximately 800~K), vibrational modes remain largely inactive, and the internal energy is governed primarily by translational and rotational contributions. Under these conditions, the gas can be treated as~\gls{CPG}, with constant specific heats and a constant ratio of specific heats, $\gamma$, more appropriately known as \textit{cold flow} in the literature. Above 800~K, vibrational excitation begins to influence thermodynamic behaviour, causing the specific heat $c_p$ and $c_v$ to vary with temperature. The gas then behaves as \emph{thermally perfect} but \emph{calorically imperfect} due to the non-constant $\gamma$. At elevated temperatures, chemical reactions become significant, initiating the dissociation and ionisation of atmospheric species within specific temperature ranges, as illustrated in \fref{figure:tempranges}~\citep{Scoggins2020}. For example, molecular oxygen (\ce{O2}) begins to dissociate at around 2500 K and is almost fully dissociated by 4000 K. These processes lead to increased thermochemical complexity in the flow.




% High-enthalpy effects have been often referred to as ``real gas” effects in the literature, which can be considered as the internal energy modes of gas absorbing, sorting, and releasing energy. These modes of energy include rotation, translation,s vibration, dissociation and exchange, electronic excitation, radiation, and ionization; therefore, total energy of the molecules or atoms are composed of the sum of these contributions and each energy mode exists at certain discrete, quantised values. A detailed description of the modes can be found in \cite[ch~11]{Anderson2006}.  

% In hypersonic flows, these modes are endothermic, acting as heat sinks absorbing heat contrary to what happens in exothermic flows like combustion. The accessibility of these modes depend on the enthalpy range of the vehicle’s flight conditions \citep{Gupta1990}. Furthermore, each mode requires different energy level before it can absorb energy, which is sometimes referred to as being activated. In order to understand the modelling implications, one must be first aware of the physical phenomena occurring to gas species at increasing temperature. In conventional flows, i.e., flowing operating at less than $800$ K, the vibrational energy mode is not fully excited, consequently, translational and rotational modes prevail dictating the total energy of the fluid. Such a gas can be assumed to be calorically perfect with a constant specific heat ratio $\gamma$, as the specific heats at constant pressure and volume respectively given as $cp$ and $cv$ remain independent of the temperature. Beyond $800$ K, the vibrational mode is excited and transforms the gas properties significantly. Considering a fluid just beyond the $800$ K mark, the vibrational energy is fully excited leading to a mixture of thermally perfect ($cp=f(T)$ and $cv=f(T)$), and calorically imperfect gas. A further increase in the temperature triggers chemical processes where the individual elements in the atmosphere starts to dissociate, further increasing the complexity of the air mixture  \citep{Fletcher,Scoggins2020}. The process for vibrational and chemical excitation within Earth's atmosphere is depicted in \fref{figure:tempranges}. 

\begin{figure}[h!]
    \centering
    \includegraphics[width=12cm]{resources/figures/introduction/schematic_tempranges.pdf}
    % \includegraphics[scale=0.8]{example-image-c}
    \caption{Temperature ranges at which vibrational excitation, dissociation and ionisation activate taken for air at 1-atm. Adapted from \cite{Anderson2019}. In general, the dissociation, and ionisation temperatures are directly proportional to the pressure.}
    \label{figure:tempranges}
\end{figure}


The emergence of additional molecular energy modes (vibrational and electronic excitation, see \fref{figure:tempranges}) necessitates a multi-temperature model to capture the relevant flow physics. At high temperatures, gas molecules store energy in distinct modes translational, rotational, vibrational, and electronic—each with different characteristic energy levels. These are often represented by separate temperatures: $T_t$ (translation), $T_r$ (rotation), $T_v$ (vibration), and $T_e$ (electronic excitation). A multi-temperature approach treats these modes individually, providing a more accurate description of the internal energy distribution. Ignoring electronic levels, a common simplification assumes the translational and rotational modes behave similarly, giving a combined translational–rotational temperature, $T_{tr}$, while vibrational modes share a single mixture temperature, $T_{v}$. This forms the basis of Park’s classic two-temperature (2T) model~\citep{Park1990}, which remains widely used for practical high-temperature flow simulations.

% These chemical and vibrational processes occur at finite time scales due to molecular collisions.  The associated time scales for each mode are detailed as: vibrational relaxation time-scale $\tau_{vib}$, chemical reaction time-scale $\tau_{chem}$ and flow time-scale (residence time) $\tau{flow}$. The coupling between each of the individual timescales dictates the distinct flow regimes, each with significant implications for the underlying flow physics. For example, they influence the rate of energy transfer between translational and internal modes, the onset and evolution of chemical reactions, and the overall thermal structure of the flow. Accurately identifying the prevailing regime is therefore essential for selecting appropriate modelling strategies and ensuring reliable predictions in high-enthalpy flow simulations. The behaviour of the flow can be characterised by the Damköhler number, defined as the ratio of the characteristic chemical or vibrational time scale to the characteristic flow time scale. For thermochemical flows, the vibrational and chemical Damköhler numbers are defined as follows:
\subsection{Flow definitions}

Chemical and vibrational processes occur over finite time scales due to molecular collisions. The relevant time scales are the vibrational relaxation time scale $\tau_{\mathrm{vib}}$, the chemical reaction time scale $\tau_{\mathrm{chem}}$, and the flow (or residence) time scale $\tau_{\mathrm{flow}}$. The relative magnitudes of these time scales give rise to distinct flow regimes, each with important implications for high-enthalpy flow behaviour. These relationships influence the rate of energy transfer between different modes, affect the onset and progression of chemical reactions, and shape the flow's overall thermal and chemical structure. The behaviour of the flow can be characterised using the Damköhler number, which compares the rate of a thermochemical process to the rate of convection. It is defined as the ratio of a characteristic chemical or vibrational time scale to the characteristic flow (or residence) time scale. For thermochemical flows, the relevant dimensionless groups are the chemical Damköhler number \gls{Dachem} and the vibrational Damköhler number \gls{Davib}, defined respectively as \citep{miromiroThesis}: \begin{equation}\label{equation:davib}
    \mathrm{Da}_{\mathrm{vib}} = \cfrac{\tau_{\mathrm{flow}}}{\tau_{\mathrm{vib}}}
\end{equation} \begin{equation}\label{equation:dachem}
    \mathrm{Da}_{\mathrm{chem}} = \cfrac{\tau_{\mathrm{flow}}}{\tau_{\mathrm{chem}}}
\end{equation} 

Depending on the relative magnitudes of the vibrational, chemical, and flow time scales, the flow regime may fall into several distinct categories. In the most general case, no assumptions are made about the relative ordering of the relevant time scales, giving rise to the \gls{TCNE} regime. All modes are of comparable magnitude, i.e., \gls{Davib} $= \mathcal{O}(1)$ and \gls{Dachem} $= \mathcal{O}(1)$. Considering the limiting case of \gls{TCE}, it is assumed that all internal modes—translational, rotational, vibrational, and chemical—adjust instantaneously relative to the flow time scales. This implies that the vibrational and chemical Damköhler numbers are much greater than one, i.e., \gls{Davib} $\gg 1$ and \gls{Dachem} $\gg 1$. For a \gls{TCF} case, vibrational and chemical modes do not have sufficient time to adjust during the flow process. As a result, the vibrational and chemical Damköhler numbers are much smaller than one, i.e., \gls{Davib} $\ll 1$ and \gls{Dachem} $\ll 1$. Isolating the Damköhler numbers allows for intermediate regimes to be examined. For example, if chemical reactions are frozen while vibrational modes adjust rapidly, then \gls{Dachem} $\ll 1$ and \gls{Davib} $\gg 1$. Conversely, if vibrational relaxation is effectively frozen while chemical reactions proceed, then \gls{Davib} $\ll 1$ and \gls{Dachem} $\gg 1$. These limiting cases provide insight into the individual roles of chemical and vibrational processes in shaping the overall flow behaviour.

% Depending on the relative magnitudes of the vibrational, chemical, and flow time scales, the flow regime may fall into several distinct categories. In the most general case, no assumptions are made about the relative ordering of the relevant time scales, giving rise to the \gls{TCNE} regime. All modes are of comparable magnitude, i.e., \gls{Davib} = $\mathcal{O}(1)$ and \gls{Dachem} = $\mathcal{O}(1)$. Considering the limiting case of ~\gls{TCE}, it is assumed that all internal modes — translational, rotational, vibrational, and chemical, equilibrate instantaneously relative to the flow time scales. This implies that the vibrational and chemical Damköhler numbers are much greater than one, i.e., \gls{Davib} $\gg 1$ and \gls{Dachem} $\gg 1$. For a \gls{TCF}, vibrational and chemical modes do not have sufficient time to equilibrate during the flow process. As a result, the vibrational and chemical Damköhler numbers are much smaller than one, i.e., $\mathrm{Da}{\mathrm{vib}} \ll 1$ and $\mathrm{Da}{\mathrm{chem}} \ll 1$. Isolating the Damköhler numbers allows for intermediate regimes to be examined. For example, if chemical reactions are frozen while vibrational modes equilibrate rapidly, then $\mathrm{Da}_{chem} \ll 1$ and $\mathrm{Da}_{vib} \gg 1$. Conversely, if vibrational relaxation is effectively frozen while chemical reactions proceed to equilibrium, then $\mathrm{Da}_{vib} \ll 1$ and $\mathrm{Da}_{chem} \gg 1$. These limiting cases provide insight into the individual roles of chemical and vibrational nonequilibrium in shaping the overall flow behaviour.



% In the most general case, \gls{TCNE}, no assumptions are made about the ordering of the time scales, which are of comparable magnitude; i.e., $\mathrm{Da}_{\mathrm{vib}} = \mathcal{O}(1)$ and $\mathrm{Da}_{\mathrm{chem}} = \mathcal{O}(1)$. When both vibrational and chemical relaxation processes are much faster than the flow time scale, the system is in \gls{TCE}, corresponding to $\mathrm{Da}_{\mathrm{vib}} \gg 1$ and $\mathrm{Da}_{\mathrm{chem}} \gg 1$. If thermal modes equilibrate quickly but chemical reactions remain too slow to respond, the flow is in \gls{TECN}, where $\mathrm{Da}_{\mathrm{vib}} \gg 1$ and $\mathrm{Da}_{\mathrm{chem}} \ll 1$. A further limiting case is \gls{TEQ}, the thermally equilibrated but chemically frozen regime, in which $\mathrm{Da}_{\mathrm{vib}} \gg 1$ and $\mathrm{Da}_{\mathrm{chem}} \ll 1$ to the extent that reactions are entirely inactive. Conversely, if the chemical processes equilibrate quickly while vibrational modes lag behind, the system is in \gls{TNCE}, defined by $\mathrm{Da}_{\mathrm{vib}} \ll 1$ and $\mathrm{Da}_{\mathrm{chem}} \gg 1$. If vibration is frozen and chemistry is also inactive, the regime becomes \gls{TFG}, with $\mathrm{Da}_{\mathrm{vib}} \ll 1$ and $\mathrm{Da}_{\mathrm{chem}} \ll 1$. In rarer cases, one may also encounter \gls{TFCE} (thermally frozen, chemically equilibrated: $\mathrm{Da}_{\mathrm{vib}} \ll 1$, $\mathrm{Da}_{\mathrm{chem}} \gg 1$), or \gls{TFCN} (thermally frozen, chemically non-equilibrium: both $\mathrm{Da}_{\mathrm{vib}}$ and $\mathrm{Da}_{\mathrm{chem}}$ are small to moderate). The most restrictive regime, \gls{TFG}, represents a fully frozen gas in which neither chemical reactions nor thermal excitation play an active role, characterised by $\mathrm{Da}_{\mathrm{vib}} \ll 1$ and $\mathrm{Da}_{\mathrm{chem}} \ll 1$. The governing formulations and flow physics change with the relative magnitudes of these timescales; therefore, accurately characterising high-enthalpy flows is crucial before addressing their interaction with high-speed phenomena.

% \newpage

\section{Boundary layer transition}


External aerodynamic flows are inherently susceptible to environmental and surface disturbances, which can destabilise the laminar boundary layer and trigger transition to turbulence. The highly fluctuating nature of turbulent flows enhances energy and momentum mixing compared to laminar flows, resulting in steeper gradients. These large energy gradients ultimately lead to significantly higher surface heat fluxes — up to eight times greater than in a fully laminar regime~\citep{VanDriest1952}. This transition is especially critical in high-speed boundary layers, which are more prone to additional disturbances that can amplify instabilities and promote earlier breakdown. The strong shockwave compressions mean that stagnation temperatures are higher, which leads to the requirement of a \gls{TPS} unit. \gls{TPS} sizing is therefore paramount to predict the transition onset and turbulent breakdown locations on the vehicle surface for the various flow regimes encountered during the mission. Consequently, the use of inaccurate transition-modelling tools leads to designs based on the worst-case scenario, with an excessive \gls{TPS} weight allocation.

% Transition is influenced by a variety of flow features, including Reynolds number, Mach number, surface geometry, inflow turbulence and acoustic noise, pressure gradients, wall roughness and porosity, boundary layer cooling or heating, and mass injection or suction through the surface \citep{Cebeci2001}. \citet{Morkovin1994} proposed a roadmap for transition in generic boundary layer flows. In general, disturbances enter the boundary layer and undergo a receptivity process that gives rise to eigenmode instabilities. These instabilities grow exponentially due to linear amplification mechanisms. The most unstable of these modes represents the primary instability and drives the initial stages of the transition process, which are subsequently followed by the development of secondary instabilities and non-linear interactions. For low-speed flows, the primary instabilities typically take the form of Tollmien–Schlichting (TS) waves, whereas for high Mach number flows, they appear as \textit{Mack modes} \citep{Mack1984}, which are two-dimensional, inviscid acoustic instabilities. A detailed description of paths to transition and the underlying mechanisms is provided in \Cref{chapter:disturbance_growth}. The main point to note here is the difference in the nature of the primary instabilities between low-speed and high-speed boundary layers, and the importance of accurately modelling these dominant modes to predict transition onset and the eventual breakdown to turbulence. The path to full breakdown strongly depends on the nature of the incoming disturbances: small-amplitude perturbations such as acoustic noise or surface vibrations typically lead to the slowest transition route, whereas surface protrusions and significant roughness can trigger bypass transition mechanisms. The present research isolates the natural transition path driven by small-amplitude disturbances to better understand the fundamental mechanisms that govern instability growth, and breakdown in high-speed boundary layers.

% Transition is influenced by a variety of flow features, including Reynolds number, Mach number, surface geometry, inflow turbulence and acoustic noise, pressure gradients, wall roughness and porosity, boundary layer cooling or heating, and mass injection or suction through the surface \citep{Cebeci2001}. \citet{Morkovin1994} proposed a roadmap for transition in generic boundary layer flows, summarising how disturbances enter the boundary layer and undergo a receptivity process that gives rise to eigenmode instabilities. These instabilities grow exponentially through linear amplification mechanisms. The most unstable of these modes represents the primary instability and drives the initial stages of the transition process, which are subsequently followed by the development of secondary instabilities and nonlinear interactions. For low-speed flows, the primary instabilities typically take the form of Tollmien–Schlichting (TS) waves, whereas for high Mach number flows, they appear as \textit{Mack modes} \citep{Mack1984}, which are two-dimensional, inviscid acoustic instabilities. A detailed description of these transition paths is provided in \Cref{chapter:disturbance_growth}. Crucially, the dominant instability mechanism differs between low-speed and high-speed boundary layers, and accurately modelling these modes is essential for predicting transition onset and breakdown to turbulence. The path to full breakdown strongly depends on the nature and amplitude of the incoming disturbances: small-amplitude perturbations, such as acoustic noise or surface vibrations, typically follow the natural transition route, whereas larger surface protrusions or significant roughness can trigger bypass mechanisms. The present study isolates the natural transition path driven by small-amplitude disturbances to clarify the fundamental processes that govern instability growth and breakdown in high-speed boundary layers.

Transition is influenced by a variety of flow features, including Reynolds number, Mach number, surface geometry, inflow turbulence and acoustic noise, pressure gradients, wall roughness and porosity, boundary layer cooling or heating, and mass injection or suction through the surface \citep{Cebeci2001}. \citet{Morkovin1994} proposed a roadmap for transition in generic boundary layer flows, summarising how disturbances enter the boundary layer and undergo a receptivity process that gives rise to eigenmode instabilities. These instabilities grow exponentially through linear amplification mechanisms. The most unstable of these modes represents the primary instability and drives the initial stages of the transition process, which are subsequently followed by the development of secondary instabilities and nonlinear interactions. A detailed description of these transition paths is provided in \Cref{chapter:disturbance_growth}. Crucially, the dominant instability mechanism differs between low-speed and high-speed boundary layers, and accurately modelling these modes is essential for predicting transition onset and the eventual breakdown to turbulence. The path to full breakdown strongly depends on the nature and amplitude of the incoming disturbances: small-amplitude perturbations, such as acoustic noise or surface vibrations, typically follow the natural transition route, whereas larger surface protrusions or significant roughness can trigger bypass mechanisms. The present study focuses on isolating the natural transition path driven by small-amplitude disturbances to clarify the fundamental processes that govern instability growth and breakdown in high-speed boundary layers.

The natural transition process has been a subject of interest since the late 19\textsuperscript{th} century, beginning with the works of Lord Rayleigh \citep{Rayleigh1878,Rayleigh1880}, who postulated the existence of sinusoidal disturbances within the incompressible boundary layer prior to transition. This mathematical framework was later expanded by \citet{Tollmien1936}, who showed that incompressible flows are stable to inviscid, temporally growing disturbances unless the two-dimensional mean velocity profile contains an inflection point. The existence of such inviscid instabilities was later confirmed experimentally by \citet{Schubauer1948}, who provided clear evidence of their downstream amplification using hot-wire measurements. When the effect of viscosity is included, these unstable waves are known as the classic Tollmien–Schlichting (TS) waves \citep{Schlichting2000}. \citet{LeesLin1946} were the first to investigate instability propagation in compressible boundary layers, although they considered only a subset of possible modes. Nonetheless, they demonstrated that a two-dimensional, inviscid compressible boundary layer remains stable unless the mean velocity and density profiles contain a generalised inflection point, analogous to the condition identified by \citet{Tollmien1936}, thus giving rise to the well-known inflection point criterion.

A detailed framework for compressible high-speed boundary layer stability analysis can be found in \citet{Mack1984} and the references therein. Mack showed that inviscid instability increases with Mach number as the generalised inflection point moves away from the wall, while viscous instability decreases and eventually disappears for two-dimensional waves at $Ma > 3$, so that viscosity has a stabilising effect in high-speed flows. It was also within this work that Mack formally classified these instability modes, distinguishing the first-mode instabilities—primarily viscous in nature and corresponding to the classical Tollmien–Schlichting waves—from the second-mode instabilities, which are essentially inviscid acoustic waves trapped within the boundary layer and become dominant at high Mach numbers. Mack performed \gls{LST} on supersonic flows and found that there are infinite number of eigenvalues due to the hyperbolic nature of the equations. The lowest order is called the second mode. As the Mach number increases, the first mode maximum growth
rate decreases. It even becomes less unstable than the second mode which becomes the
dominant instability mode for M > 4. In low speed flows, a 2D wave is most unstable
whereas an oblique first mode wave is most unstable in a supersonic boundary layer (M
< 4). In a hypersonic boundary layer, the 2D second-mode waves grow faster. Frequencies of first-mode waves are around tens of kilohertz while second mode waves can be found around hundreds of kilohertz. Both LST and experiments showed that wall cooling had the effect of stabilizing the first mode and destabilizing the second mode. The second mode could even be more unstable than the first mode on a cooled wall at smaller Mach numbers. Numerical simulations showed that the second mode could be stabilized by suction and favorable pressure gradients.

Unlike the Tollmien-Schlichting waves mentioned above which can be stabilized by
wall cooling, the acoustic modes are destabilized, an eect which is mainly caused by
the reduction in boundary layer thickness that occurs when the wall is cooled. This
destabilization can be thought of in the following way: as the boundary layer becomes
thinner, the wavelength of the second mode waves also decreases since they are reverberating within the boundary layer. The amplication per unit wavelength remains
approximately constant, and hence the amplication per unit distance (the growth
rate) increases as the boundary layer is made thinner, since there are more waves per
unit distance. This explanation is demonstrated quantitatively in Section 5.4.2.

For boundary layer fows that have a high stagnation enthalpy, the temperature
in the interior of the boundary layer can become large enough that the vibrational
energy of polyatomic molecules becomes signicant, such that the gas cannot
be treated as calorically perfect. Dissociation reactions can also become signi-
cant at suciently high enthalpies. Furthermore, the rates of chemical reactions or
vibrational-translational energy transfer can be comparable to 
ow timescales, leading to a nonequilibrium ow. A number of researchers have considered these nonequilib turbations appears in the 
ow, for instance, behind a small, sharp-edged roughness
element. The vertical velocity generated by these vortices then carries low momentum uid up from the wall and high momentum 
uid down from the freestream, thereby producing streaks of high and low streamwise velocity. The existence of this mechanism
has been demonstrated in experiments involving roughness elements placed in
incompressible boundary layers. Such experiments have reproduced many features
predicted by transient growth calculations, including the downstream evolution of
disturbance kinetic energy and the shape of the amplied disturbances (White, 2002,
White et al., 2005, Ergin and White, 2006).




% In the most general case, \gls{TCNE}, no assumptions are made about the ordering of the time scales which are of comparable magnitude, i.e., $\mathrm{Da}_{\mathrm{vib}} =\mathcal{O}(1)$ and $\mathrm{Da}_{\mathrm{chem}} =\mathcal{O}(1)$. 

% When vibrational relaxation is much faster than the flow time scale, and both are faster than the chemical reactions, the regime is described as chemically frozen and in \gls{TEQ}, where $\tau_{\mathrm{vib}} \ll \tau_{\mathrm{flow}} \ll \tau_{\mathrm{chem}}$. If both vibrational and chemical processes relax much more rapidly than the flow, i.e., $\tau_{\mathrm{chem}} \sim \tau_{\mathrm{vib}} \ll \tau_{\mathrm{flow}}$, the flow is in \gls{TCE}. On the other hand, when vibrational relaxation is much slower than both flow motion and chemical activity, i.e., $\tau_{\mathrm{vib}} \gg \tau_{\mathrm{flow}} \sim \tau_{\mathrm{chem}}$, the flow is in a thermal equilibrium but chemical non-equilibrium (TEQ–CNE) state. Conversely, if chemical relaxation is fast compared to flow and vibrational processes, i.e., $\tau_{\mathrm{chem}} \ll \tau_{\mathrm{flow}} \sim \tau_{\mathrm{vib}}$, the flow exhibits thermal non-equilibrium with chemical equilibrium (TNE–CE). These classifications are essential for guiding the appropriate choice of physical models in high-enthalpy flow simulations.

% Depending on the relative magnitudes of these time scales, the flow regime can be classified into the following scenarios: 

% \begin{itemize}
%     \item \textbf{Thermal and chemical non-equilibrium} TCNE is the most general case, with no assumptions about the relation between the time scales, i.e. $\tau{vib} \sim \tau{chem} \sim \tau{flow}$. 
%     \item \textbf{Chemically frozen and thermal equilibrium} CF-TEQ occurs if the vibrational relaxation time occurs substantially faster than the flow times scale and both are smaller than the chemical timescale $\tau{vib} \ll \tau{flow} \ll \tau{chem}$.
%     \item \textbf{Thermal and chemical equilibrium} TCEQ occurs when both the chemical and vibrational time scales are significantly faster than the flow time scale: $\tau{chem} \sim \tau{vib} \ll \tau{flow}$
    
%     \item \textbf{Thermal equilibrium, chemical non-equilibrium} TEQ-CNE is assumed when the chemical activity and the flow motion are significantly faster than the vibrational relaxation: $\tau{vib} \ll \tau{flow} \sim \tau{chem}$. 
    
%     % The two-temperature assumption (\eref{equation:twotemperature}) could be reduced down to a one temperature model when this condition occurs. 
%     \item \textbf{Thermal non-equilibrium, chemical equilibrium} TNE-CE, where the chemical relaxation is significantly slower than the chemical process and flow motion $\tau{chem} \ll \tau{flow} \sim \tau{vib}$, the system is assumed to be in a thermally non-equilibrium state. 
% \end{itemize}



% \section{Laminar turbulent transitions}
% % \section{Transitional flows}
% The transition from a laminar to a turbulent boundary layer is a critical aerodynamic phenomenon with significant implications for skin friction, heat transfer, and aerodynamic drag. The destabilisation of an otherwise laminar flow into a chaotic, turbulent regime had attracted interest since the late 1800s when Osborne Reynolds first demonstrated that laminar flow could become unstable under certain conditions. Over the decades, extensive research has sought to understand, predict, and control this process through theoretical and experimental efforts. Within this broader context, \citet{Morkovin1994} laid out a conceptual foundation that synthesised these developments into the modern view of the multiple paths to transition that can occur in boundary layer flows. \fref{figure:transitionmarkovin}, reproduced from \citet{Fedorov2011}, presents a schematic overview of these paths to transition, which applies to both low-speed and high-speed boundary layers, although the specific instability modes involved differ between the two regimes.


% \begin{figure}[t!]
%     \centering
%     \fontsize{9}{10}\selectfont\includesvg[width=0.50\textwidth]{resources/figures/introduction/schematic_fedorov.svg}
%     \caption{Diagram of pathways by which a boundary layer might transition to turbulence \citep{Fedorov2011}.}
%     \label{figure:transitionmarkovin}
% \end{figure}

% % The transition begins with the appearance and subsequent amplification of disturbances, which can arise from environmental sources such as freestream turbulence, acoustic noise, surface roughness, or shock interactions, depending on the flow regime. These environmental disturbances interact with the boundary layer, appearing as fluctuations in the base flow that generate oscillatory modes whose amplitude, frequency, and phase set the initial conditions for the instability waves that develop and ultimately lead to breakdown. This process, known as the receptivity mechanism, converts external disturbances into instability waves; their initial amplitude plays a key role in determining how and when the transition occurs. Because environmental disturbances are often complex and contain a wide range of information, the receptivity problem remains largely unsolved for compressible boundary layers, although several studies have tried to tackle the issue at hand. 

% % Transition begins with the appearance and subsequent amplification of disturbances originating from environmental sources such as freestream turbulence, acoustic noise, surface roughness, or shock interactions, depending on the flow regime. These external disturbances interact with the boundary layer, manifesting as fluctuations in the base flow that generate oscillatory modes whose amplitude, frequency, and phase define the initial conditions for instability waves that develop and eventually breakdown. This process, known as receptivity, converts external disturbances into instability waves whose initial amplitude plays a critical role in determining the onset and location of transition. In hypersonic boundary layers, key destabilising factors influencing transition in hypersonic flows include: \textbf{Freestream disturbances}: Acoustic waves, entropy fluctuations, and vorticity perturbations in the freestream can penetrate the boundary layer and trigger instability growth. \textbf{Surface roughness}: Imperfections on the surface can introduce localised perturbations that bypass classical instability mechanisms, leading to premature transition. \textbf{Shock interactions}: Shock-induced perturbations can introduce unsteady pressure and velocity fluctuations, modifying the growth of boundary layer instabilities. \textbf{High-temperature effects}: Vibrational excitation, dissociation, and recombination reactions in high-enthalpy flows alter the stability properties of the boundary layer, often enhancing instability growth. Because environmental disturbances are typically complex and broadband, the receptivity problem remains largely unsolved for compressible boundary layers, despite numerous efforts to address it.

% Transition begins with the appearance and amplification of disturbances from environmental sources such as freestream turbulence, acoustic noise, surface roughness, or shock interactions. These interact with the boundary layer, generating oscillatory modes whose amplitude, frequency, and phase define the initial conditions for instability waves that grow and ultimately breakdown. This \textit{receptivity process} governs how external disturbances are converted into instability waves, strongly influencing when and where transition occurs \citep{Morkovin1994}. In hypersonic boundary layers, key destabilising factors include freestream perturbations (acoustic, entropy, and vorticity fluctuations), surface roughness that introduces localised bypass routes, shock interactions that create unsteady pressure and velocity fields, and high-temperature effects like vibrational excitation and dissociation, all of which can significantly enhance instability growth \citep{vandereynde2015,bitter2015}. Because these environmental disturbances are inherently complex and broadband, the receptivity problem remains largely unresolved for compressible boundary layers, despite extensive research \citep{dCerminara2017} Ma Zhong.

% In the most general case, low-amplitude oscillatory modes are triggered within the boundary layer and amplify through linear modal growth, followed by parametric instabilities and mode interactions that transfer energy between modes and lead to breakdown to turbulence; this transition pathway is depicted in \fref{figure:transitionmarkovin} as \textit{path A}. This path has been widely studied since the pioneering work of Tollmien and Schlichting, which established the well-known two-dimensional Tollmien--Schlichting instability waves. For higher disturbance levels, the superposition of perturbations may first exhibit algebraic transient growth before feeding into modal amplification (\textit{path B}). As the disturbance level increases further, transient growth can lead directly to parametric instabilities and mode interactions or, if sufficiently strong, bypass these stages entirely and transition straight to breakdown (\textit{paths C} and \textit{D}). Only \textit{paths A, B,} and \textit{C} are relevant for external hypersonic boundary layers, where disturbances grow through receptivity, transient growth, and modal amplification. \textit{Path D} mainly applies to internal flows with elevated turbulence levels, while \textit{path E} represents extreme forcing such as severe surface roughness or intense freestream turbulence where linear growth mechanisms are completely bypassed.


% Most of the groundwork done for  

% Ok, hear me out, most of the groundwork done for transition prediction was on incompressible flows which is not completely true but we say that anyway. The first studies about boundary layer stability was performed my Lees and Lin, who were quantifying the inflection point phenomenon, then came the Tollmein-Schliting waves, confirmed by further experiments by Schuberger or whatever their names are. It was not until Mack studies the stability of supersonic boudnary layer that we have a rigiorous definition of boundary layers and what they actually are. Mack performed stability analysis on supersonic base flows and characterised instability modes into numbers (first, second and third mode, so on). A detailed description of these instabilities is discussed in chapter 4, but the general gist of it is, first modes are promiment in subsonic incompressible flows, often associated with lower frequencies and the second mode is an acoustic trapped wave that is a feature of flows past supersonic regime. 



% The primary focus of the current thesis is to understand the modal decomposition of \textit{Path A}, 

% The most general transition process driven by freestream disturbances, known as Path A, involves low-amplitude disturbances that undergo exponential growth of eigenmodes. This linear amplification leads to parametric instabilities and mode interactions, which eventually cause breakdown to turbulence.  

% The most general transition process driven by freestream disturbances involves low-amplitude perturbations that undergo exponential growth of eigenmodes. This linear amplification leads to parametric instabilities and mode interactions, which effectively transfer energy between modes and eventually cause breakdown to turbulence. This pathway is illustrated in \fref{figure:transitionmarkovin} as \textit{path a}.

% In the most general case, low-amplitude oscillatory modes are triggered within the boundary layer and amplify through linear modal growth, followed by parametric instabilities and mode interactions that transfer energy between modes and lead to breakdown to turbulence. This transition pathway is depicted in \fref{figure:transitionmarkovin} as \textit{path a}. For slightly higher disturbance levels, the superposition of perturbations may exhibit algebraic transient growth before transitioning to modal growth; this transient growth mechanism, \textit{path b} which ultimately feeds into the same eigenmode amplification process as \textit{path a}, but is preceded by an initial phase of non-modal growth. For even higher disturbance levels, the flow may undergo transient growth that either leads directly to parametric instabilities and mode interactions (path C) or, if the disturbance level is sufficiently high, bypasses these stages entirely and transitions straight to breakdown (path D). Only paths A, B, and C are relevant for studying external hypersonic boundary layers, where disturbances grow through receptivity, transient growth, and modal amplification. Path D mainly applies to internal flows with elevated turbulence levels, where modal growth is bypassed altogether. Path E represents extreme disturbance environments — such as severe surface roughness or very high freestream turbulence — where linear mechanisms are entirely bypassed, leading directly to breakdown.




% Only paths A, B and C are
% sensible for the study of external hypersonic flows. Path D is primarily for internal flows
% at an elevated turbulence level. Path E represents the case of very large amplitude forcing
% (roughness and high freestream turbulence) where the growth of linear disturbances is
% bypassed




% The transition from a laminar to a turbulent boundary layer is a critical phenomenon in aerodynamics, with significant implications for skin friction, heat transfer, and aerodynamic drag. This process is initiated by amplifying flow disturbances, which may originate from various environmental sources such as freestream turbulence, acoustic noise, surface roughness, or shock interactions. The specific pathway to turbulence depends on the boundary layer's stability characteristics and the initial perturbations' nature and amplitude. A schematic overview of the various transition routes relevant to a flat plate boundary layer is shown in \fref{figure:transitionmarkovin} reproduced from \citep{Fedorov2011}.

% The transition from a laminar to turbulent boundary layer is a critical phenomenon in aerodynamics, with significant implications for skin friction, heat transfer


% \begin{figure}[h!]
%     \centering
%     % \includegraphics[width=13cm]{resources/figures/introduction/schematictempranges.png}
%     \includegraphics[scale=0.8]{example-image-c}
%     \caption{Markovin transition path to turbulence.}
%     \label{figure:transitionmarkovin}
% \end{figure}

% \begin{figure}[h!]
%     \centering
%     % \includegraphics[width=13cm]{resources/figures/introduction/schematictempranges.png}
%     \includegraphics[scale=0.8]{example-image-c}
%     \caption{Markovin transition path to turbulence.}
%     \label{figure:transitionmarkovin}
% \end{figure}





% According to the framework established by \cite{Morkovin}

% In incompressible and low-speed flows, the transition is typically governed by the amplification of two-dimensional Tollmien–Schlichting (TS) waves, which correspond to the \textit{first-mode} instability in Mack’s classification \citep{Mack1990}. These disturbances undergo exponential growth, eventually leading to nonlinear flow features such as $\Lambda$-vortices, hairpin vortices, and turbulent spots, culminating in fully developed turbulence. In contrast, hypersonic boundary layers are subject to additional instability mechanisms arising from compressibility effects, high-temperature-induced chemistry, and shock interactions. Key destabilising factors influencing transition in hypersonic flows include: \textbf{Freestream disturbances}: Acoustic waves, entropy fluctuations, and vorticity perturbations in the freestream can penetrate the boundary layer and trigger instability growth. \textbf{Surface roughness}: Imperfections on the surface can introduce localised perturbations that bypass classical instability mechanisms, leading to premature transition. \textbf{Shock interactions}: Shock-induced perturbations can introduce unsteady pressure and velocity fluctuations, modifying the growth of boundary layer instabilities. \textbf{High-temperature effects}: Vibrational excitation, dissociation, and recombination reactions in high-enthalpy flows alter the stability properties of the boundary layer, often enhancing instability growth.  


% A defining feature of hypersonic flows is the presence of \textbf{Mack modes}, or \textit{second-mode} instabilities, which arise from the interaction of trapped acoustic waves between the wall and the sonic line. The first theoretical studies of these compressible eigenmodes were conducted by Mack, who extended the classical linear stability theory (LST) developed by Orr and Sommerfeld. The second-mode instability in hypersonic boundary layers was first investigated numerically by \citet{Malik1990a}, who demonstrated its rapid amplification. Recent studies, such of \cite{Bitter2015} and \citet{miromiro2019}, have shown that dissociation effects in high-enthalpy regimes tend to destabilise second-mode instabilities, increasing their amplification rates. These findings also highlight the strong sensitivity of flow stability to variations in base flow conditions and wall properties.

% Building on these insights, recent advances in boundary-layer stability analysis have focused on computing second-mode $N$-factor envelopes using linear stability theory (LST), typically assuming frozen base flows \cite[p.~12-7]{miromiroThesis}. However, in post-shock regions—particularly near stagnation points—the boundary layer often approaches thermochemical equilibrium conditions \citep{Passiatore2021}, which challenges the validity of the frozen-flow assumption. This highlights the need to investigate how thermochemical state and wall boundary conditions influence stability characteristics. A deeper understanding of these effects is crucial for improving transition prediction models and informing control strategies in high-speed vehicle design.

% \section{Stability studies}
% Mention past stability studies that have gone into understanding the basic hypersonic regime

% \section{Numerical investigations}
% Talk about both temporal and spatial simulations for high enthalpy boundary layers


% \newpage

\section{Recent Progress in Hypersonics}\label{section:introduction/recentprogress}

Recent progress in hypersonic aerothermodynamics has focused on improving the understanding and prediction of complex flow phenomena, including boundary layer transition, shock–boundary layer interactions, and thermochemical nonequilibrium effects.




The primary motivation for this work is to advance the understanding of the different paths to transition in boundary-layer flows and to bridge the gap between the well-understood linear stages and the more complex nonlinear breakdown to turbulence. Accurate prediction of transition is critical for reliable estimates of aerodynamic drag, heat transfer, and vehicle performance, yet experimental investigations are often constrained by the energy demands and practical limitations of recreating realistic conditions, especially for high-speed flows. \gls{DNS} provides a powerful alternative by delivering high-fidelity, fully resolved data that can capture the entire transition process under controlled conditions. By using DNS, this work aims not only to clarify how freestream disturbances and receptivity mechanisms feed into the transition process, but also to generate benchmark data that can guide improved models and future experimental efforts.

% \subsection{Instability growth}
Considering a generic leading edge structure shown in \fref{figure:hypersonicwedge}, the transition derives from the flow conditions and any receptivity in the leading edge region, therefore, understanding the laminar boundary layer behaviour adjacent to the stagnation-point is important to understand the potential growth to turbulence transition. Understanding of the stagnation point flow in complete thermochemical non-equilibrium is
important since this region states the conditions for the start of transition.


% \begin{figure}[h!]
%     \centering
%     \includegraphics[scale=0.36]{resources/figures/introduction/schematicshperecone.png}
%     \caption{Schematic of a sphere-cone geometry in a hypersonic flow field showcasing various regions in frozen, equilibrium and non-equilibrium conditions adapted from \citep{williams2021locally}.}
%     \label{figure:hypersonicwedge}
% \end{figure}

% The stagnation-point flow around a blunt body has been a subject of interest for since Hiemenz's numerical studies in 1911. The widely used approach for a dissociating gas in compressible flow was developed by \cite{FayRiddell1958}. The method provides an exactly self-similar solution to the Naiver-Stokes equations in dissociating gas to understand the stagnation point behaviour in the vicinity of a blunt geometry under various wall catalysis conditions. However, it should be mentioned that several authors (e.g. \citep{Lee2022} and \citep{Rocher2022}) have shown that these equations contain several approximations in their crude estimation of several crucial properties. 

% From an experimental standpoint, several empirical equations have been developed over the years to understand the heat-flux at the stagnation-point, most notably, Detra-Hidalgo \citep{Detra1961}, Detra-Kemp-Riddell, Brandis-Johnson \citep{Brandis2014} and Verant-Sagnier \citep{Sourgen2015}. More recently, \cite{Rocher2022} have extended the Fay-Riddell framework, incorporating semi analytical methods to include thermochemical equilibrium, while determining the surface species concentration in the stagnation point. However, the empirical nature of these correlations maybe inadequate for conditions outside the tested experimental regimes, prompting the requirement for more focused experiments to accurately capture the physics of the flow \citep{donaldson2021a}.  
% Concluding from the studies above, a sufficient self-similar framework incorporating chemical and vibrational non-equilibrium in various flow regions has been a vastly unexplored topic, providing scope for future research. 

The standard experimental ground-based hypersonic facilities rarely operate at the conditions similar to that of a low altitude high-speed vehicle, as the facility power requirements to replicate the exact combination of velocity, flow density and therefore the stagnation enthalpy can be overwhelming to the current electrical grids \citep{passiatore2022}. Laminar and turbulent flows have received significant attention in past albeit at sub-orbital enthalpies, with approximations about recombination/dissociation reaction mechanisms. Several of these studies focusing on simplified turbulent hypersonic boundary layers with minimal chemical reactions were developed in the past in \cite{martin1998} , although it was not until \cite{duanmartin2011} that a rigorous differentiation between low-enthalpy and high-enthalpy calculations was provided, for a gas with little chemical activity neglecting thermal non-equilibrium. The studies have mostly considered temporal evolution of the boundary layer with little or no chemical activity. Therefore, little attention was directed towards understanding the nature at the extreme conditions of hypersonic flight \citep{Candler2019}.

% (Martin \& Candler, 1998, 1999, 2001) + Dyab \& Martin (2009)
% dCerminara2017,Cerminara2019,cerminarasandham2020

Transition to turbulence in hypersonic boundary layers has been studied in the past, both in \textit{`cold flow'} and high-enthalpy regime at suborbital enthalpies, despite the challenges imposed by the later environment. The studies concern stages of transition through perturbations in the laminar base flow through linear and non-linear stability analyses. \cite{Cerminara2019, cerminarasandham2020} studied the receptivity mechanisms at the leading edge of a blunt body in cold hypersonic flow without considering finite-rate chemical models. Studies of transition and control for chemically perfect gases, and flows in chemical equilibrium were first performed by  \citep{Malik1989, Malik1990b, Malik1990a, Malik1990c}. This was extended further by \cite{Stuckert1994} and \cite{Hudson1997}, who introduced chemical and vibrational non-equilibrium extensions with Park's two-temperature model. The studies conclusively show that high enthalpy effects bring non-negligible contribution to boundary layer stability; the endothermic nature of the reactions remove energy from the thermal boundary layer, leaving the flow more subjected to second-mode (Mack modes) disturbances, almost mimicking wall-cooling effects. \cite{Malik2003} and \cite{Johnson2005}, conducted parabolised stability equations (PSE) to the subject, with the former study concluding that some dissociation reactions actually reduce the growth of perturbations caused by the chemical reactions on the base flow. \cite{Johnson2005} has also performed simulations using multi-species models to better capture the stability during transition. \cite{Franko2010} concluded that the chemical and viscosity model used in the analysis has a substantial effect on the stability characteristic. Consequently, \cite{Klentzman2013} showed that viscosity stabilises the base flow while investigating viscous effects on the receptivity of boundary layers and \cite{miromiro2018} extended the work on an adiabatic flat plate, observing the flow to most destabilise when mass diffusion is hundreds of times faster than viscous diffusion and further investigated the effects of choice of thermochemical and thermophysical framework on the transition prediction \citep{miromiro2019}. Furthermore, work of \cite{wartmann2018} show that thermochemical non-equilibrium approach for the mean flow and the stability code are almost identical to the approach of considering just the base flow in thermochemical non-equilibrium. Most of the work in the literature focused on the modal growth of the transition process predominantly for sub-orbital enthalpies. The stability of high-speed flows at low altitudes, where the Reynolds numbers are significantly higher and a comprehensive study of the evolution of the boundary layer to complete turbulence transition in such flows has been largely unexplored areas.  

\subsection{Experiments and Modelling}
\gls{DNS} of flow close to such extreme hypersonic regimes with thermochemical framework has recently appeared in the literature, with the first studies performed by \cite{marxen2013} on laminar boundary layers with instabilities suitable for hypersonic cruise conditions, whilst an extension study \citep{marxeniaccarinomagin2014} focused on non linear instability in the presence of finite rate chemistry without considering thermal non-equilibrium, concluding the dominant influence of chemical properties on primary disturbance amplitude, neglecting secondary instabilities during transition. Simulations of chemical non-equilibrium have received far greater attention as opposed to thermal non-equilibrium, with the Mutation++ library \cite{Scoggins2020}, providing an extensive state of the art computation ability of physiochemical properties and has been widely used in the chemical studies concerning hypersonic flows. \cite{DiRenzo2020} introduced their HTR finite-difference solver for hypersonic aerothermodynamics for exascale computing, \cite{direnzourzay2021} used the routine to produce spatial evolution of such boundary layers to fully turbulent state. \cite{Passiatore2021} studied the complete spatially evolving turbulence transition process in thermochemical non-equilibrium for a quasi-adiabatic non-catalytic wall, where a second-mode transition was triggered with a strong amplitude and observed chemical activity to have little effect on the sub-harmonic transition. These studies were further extended to understand transition process in flat-plate boundary layer configurations with non-zero pressure gradients, i.e. oblique shock interactions by \citep{Volpiani2021} using another finite-difference code \texttt{CHARLIE}, and \cite{passiatore2023}, which is essentially an thermochemical extension of the shock wave interaction work of \cite{sandham2014} in hypersonic cold flow. The studies have concluded that mean velocities and turbulent intensity profiles are similar to those of cold hypersonic flows, in fact, the presence of varying species mass fractions caused a reduction in the turbulent fluctuations. At this point,  a complete understanding of chemical and thermal influences on modal stability for high Reynolds number, high Mach number flows is available in the literature. 

\section{Motivation}

\subsection{Aims and objectives}\label{subsection:aimsandobjectives}

Following the aforementioned motivation, the project intends to develop a fundamental understanding of the mathematical and computational framework for high-enthalpy hypersonic boundary layer in thermochemical non-equilibrium and a physical understanding of transition to turbulence where ionisation will not be considered. Initially, the research intends to provide an extension of self-similar boundary layer theory which has received little attention in the low-altitude high-enthalpy region with various species and thermochemical non-equilibrium effects. Most of the linear-stability theory studies in the literature have considered a freestream frozen base flow \cite{miromiroThesis}, however, post shock flow near stagnation point region are predominantly in equilibrium conditions which are not well understood. Furthermore, the project intends to understand the transition mechanism in high-Mach number, high-Reynolds number regions under different wall conditions. Part of the work is therefore dedicated to extend the current in house finite difference compressible solver OpenSBLI to include thermal and chemical no equilibrium with appropriate models and shock capturing techniques. Furthermore, the project intends to understand the mechanisms of transition to turbulence in presence of extreme non-linear effects. Taking into account the current state of the art, following objectives are proposed to address the aforementioned crucial gap in knowledge of transitions on a wedge and the leading edge of a blunt body in thermochemical non-equilibrium. 




% Furthermore, the project intends to understand the bypass transition in the realms of hypersonic flow under different wall conditions. In particular, the aims to understand the nature of the flow in transition region to turbulence in presence of extreme non-linear effects. 
% Part of the work is dedicated to extend the current in-house compressible solver to include thermal and chemical nonequilibrium with appropriate models and shock chapturing techniques. 


% \begin{enumerate}[topsep=0pt,itemsep=-1ex,partopsep=1.5ex,parsep=3.5ex]
%     \item To extend the analytical set of Fay-Riddell equations to include multiple species,  thermochemical non-equilibrium, and Park's two-temperature model.
%     % \item To perform DNS simulations of a flow around a cylinder by extending OpenSBLI for dissociating gas incorporating equations \sref{section:literature/navier-stokes} for necessary validations. 
%     \item Extend the in-house compressible solver OpenSBLI to include thermochemical framework with chemical and vibrational non-equilibrium, including validation for flow around a cylinder to understand the stagnation point behaviour. 
    
%     % \item To perform DNS simulations of flow around a cylinder to understand the stagnation point behaviour under equilibrium freestream, including validation for .
%     % \item To accurate incorporate DNS methods with suitable shock-capturing scheme to simulate the laminar-turbulent transitions on a flat plate. 
    
%     % \item To accurate incorporate DNS methods with suitable shock-capturing scheme to simulate the laminar-turbulent transition. 
%     % \item 
%     \item To investigate the effects of LST and its application in e\textsuperscript{N} method to accurately predict the behaviour of instability modes and their sensitivities to the high-enthalpy base flow inaccuracies.

    
%     \item Understand the mechanics of temporally and spatially developing turbulent flow region, using DNS of spreading turbulent spots and wedge of developing turbulent flow. 

%     % \item 
%     % \item Apply e\textsuperscript{N} method to the modal analysis to understand the sensitives of turbulence transition.
%     % \item 
% \end{enumerate}

% The physics of hypersonic flows differs vastly from the subsonic and supersonic counterparts. Most vehicles in hypersonic flows encounter strong shock-wave compression resulting in high-enthalpy gas environment that almost always associates with non-equilibrium thermodynamic and  chemical phenomena. These high-enthalpy conditions often render the calorically perfect gas models moot, consequently a multi-species model to capture the chemically reacting atmosphere is required. Further complexities such as shock-waves, in addition of these chemical reactions, add more challenges to accurately predict the loads on the vehicle.  

% still a foreigner concept along with the understanding of the mechanisms to transition, despite the great deal of attention received by them over the years. 

% The outline of the current report is as follows: necessary theoretical framework for the project is suggested in Chapter 2, whilst the numerical framework along with the simulation method is ascribed in Chapter 3. Chapter 4 covers the self-similar studies for both stagnation region and numerical inflow conditions with thermochemistry. Currently extensions to the compressible Navier-Stokes DNS solver to include thermochemistry are mentioned in Chapter 5 along with studies on cylinder flow, flat plate simulations and preliminary transition results. Finally, conclusions of the current work are drawn along with necessary framework for the thesis completion. 

\newpage

\section{Thesis outline}
The outline of the report is as follows: necessary theoretical framework 




The outline of the report is as follows: necessary theoretical framework for the project is suggested in \Cref{chapter:numerical_methods}, and the numerical framework employed for the theoretical framework along with the simulation method is ascribed in Chapter 3; which also includes the implementation of the catalytic wall boundary condition on a flow around a cylinder. Chapter 4 covers the self-similar studies for stagnation point flow past a dissociated shock layer, along with suggestions for inflow boundary conditions for compressible thermochemical simulations. Current endeavours to extend the compressible Navier-Stokes DNS solver OpenSBLI to include thermochemistry is mentioned in Chapter 5, which provides preliminary studies on flat plate laminar and transitional boundary layers in both cold and high-enthalpy supersonic flows. Finally, conclusions of the current work are drawn in Chapter 6 along with necessary future framework for the thesis completion. 